% CLAUDE REPORT::B3-EQUIVALENCE-0001::2026-11-05
\documentclass[11pt]{article}
\usepackage[margin=1in]{geometry}
\usepackage{amsmath,amssymb,amsthm}
\newtheorem{proposition}{Proposition}
\title{CLAUDE B3-EQUIVALENCE-0001: what the bridge B3 costs for a fixed system}
\date{2026-11-05}
\begin{document}\maketitle
\noindent\textbf{STATUS:} elementary logical observations (PROVED/EXACT) plus KNOWN THEOREM citations. No new lower bound. \textbf{P versus NP remains OPEN.}

\section*{Corpus reviewed / duplication audit}
Read: worst-case refocus (2026-11-02\_1204), LSPV-1/336 (size-optimality $\Rightarrow$ B1,B2,B3), LSPV-1/395 (uncomputable node profile), AGENTS.md. Proposition 1 is a restatement overlapping LSPV-1/334/336; Proposition 2 is a precise packaging of known material. Nothing here changes root \texttt{main.tex}; no specialist region was edited.

\section*{Statements}
For a fixed Cook--Reckhow system $R$ let B3$(R)$: ``$\mathrm{SAT}\in\mathrm P\Rightarrow\overline W_R(N)\le\mathrm{poly}(N)$''.

\begin{proposition}[B3 is only as cheap as the lower bound is unknown]
(a) If $R$ is polynomially bounded, B3$(R)$ holds trivially. (b) If $\overline W_R$ is \emph{unconditionally} superpolynomial (e.g.\ Resolution, Haken 1985; bounded-depth Frege), then B3$(R)\iff \mathrm P\ne\mathrm{NP}$. (c) Hence for any $R$ with a known superpolynomial lower bound, B3 cannot be proved independently of $\mathrm P\neq\mathrm{NP}$, and a lower bound plus B3 is circular. (d) For $R$ with no known lower bound, B3$(R)$ combined with a lower bound proves $\mathrm P\ne\mathrm{NP}$; B3 holds if $R$ is optimal (LSPV-1/336), but then a superpolynomial bound for $R$ is equivalent to $\mathrm{NP}\ne\mathrm{coNP}$, a strictly stronger target than $\mathrm P\ne\mathrm{NP}$.
\end{proposition}
\begin{proof}
(b): the consequent is false, so the implication is equivalent to the negation of the antecedent. (d): Cook--Reckhow: a polynomially bounded system exists iff $\mathrm{NP}=\mathrm{coNP}$; an optimal $R$ polynomially simulates it. \qed
\end{proof}

\begin{proposition}[Extended Frege conditional, KNOWN THEOREM packaging]
If SAT is decided by a polynomial-time $A$ and the soundness of $A$ is provable in $\mathrm{PV}$ (equivalently $S^1_2$ at the $\forall\Sigma^b_1$ level), then by Cook's translation EF has polynomial-size proofs of all tautologies. Without a provable correctness proof for $A$ no conclusion about EF follows (Krajíček--Pudlák 1989). So B3(EF) is reduced to a reflection/formalisation question about a hypothetical $A$, not proved.
\end{proposition}

\section*{Consequence for the program}
The only non-circular route is a bridge for a system $R$ chosen so that (i) a superpolynomial lower bound is provable and (ii) $R$ polynomially simulates every $Q_A$. By (b) and (d) these two requirements pull in opposite directions: (i) known $\Rightarrow$ B3 equals the goal; (ii) $\Rightarrow$ the lower bound is as hard as $\mathrm{NP}\ne\mathrm{coNP}$ for optimal systems. Escape would need a system with a \emph{restricted} simulation property (only against deterministic polytime $A$ under a formalisable correctness hypothesis) and a lower bound technique avoiding the known barriers; neither is supplied.

\section*{Literature check (2026-11-05)}
Searched p-optimal / optimal systems / EF. Found nothing new bearing on B3: optimal-system existence and EF lower bounds remain open (Krajíček survey); Alekseev (Theory of Computing 22, 2026) gives an exponential lower bound for Res-Lin/polynomial calculus with extension, a weaker system than EF. Classification: failed transfer for B3.

\section*{Next obligations}
Red team: find an $R$ with provable lower bound that is nonetheless simulation-closed under a nontrivial class of $A$ (would falsify the tension above). Auditor: charge verification cost in time/space for $Q_A$.

% END_OF_VERIFIED_REPORT::CLAUDE::B3-EQUIVALENCE-0001
\end{document}
